% \section{Rancangan}

% \subsection{Arsitektur Sistem}

Arsitektur sistem audit trail yang diusulkan dapat dilihat pada Gambar~\ref{fig:arsitektur-sistem}. Sistem yang dirancang merupakan sebuah \textit{Audit Trail Service} yang diintegrasikan ke dalam arsitektur DecMed untuk menyediakan mekanisme pembangkitan, penyimpanan, dan pengambilan kembali (\textit{retrieval}) \textit{audit trail} sebagai bukti digital. Berbeda dengan komponen aplikasi pada DecMed yang berfokus pada penyediaan layanan Rekam Medis Elektronik (RME), sistem audit trail berfungsi sebagai infrastruktur pendukung yang mencatat aktivitas-aktivitas penting yang terjadi pada sistem sehingga dapat digunakan sebagai bukti dalam proses audit maupun investigasi.



Secara umum, arsitektur sistem terdiri atas empat bagian utama, yaitu komponen DecMed sebagai penghasil \textit{audit event}, \textit{Audit Trail Service}, \textit{Audit Repository}, dan auditor sebagai pihak yang melakukan pengambilan \textit{audit trail}. Komponen DecMed meliputi seluruh layanan dan aplikasi yang telah diidentifikasi pada tahap analisis sebagai sumber \textit{audit event}, seperti aplikasi klien, \textit{proxy re-encryption server}, \textit{smart contract}, maupun komponen lainnya. Setiap aktivitas yang memenuhi kriteria \textit{audit event} akan dikirimkan menuju \textit{Audit Trail Service} untuk diproses.

\textit{Audit Trail Service} merupakan komponen utama yang dikembangkan pada penelitian ini. Komponen ini bertanggung jawab untuk menerima \textit{audit event}, membangkitkan \textit{audit record} sesuai dengan skema yang telah ditentukan, serta menerapkan mekanisme yang diperlukan untuk menjaga integritas, keaslian (\textit{authenticity}), dan \textit{non-repudiation} dari setiap \textit{audit record}. Setelah diproses, \textit{audit record} akan diteruskan menuju repositori audit untuk disimpan sesuai dengan kebijakan retensi yang berlaku.

Penyimpanan \textit{audit trail} dilakukan pada \textit{Audit Repository} yang dirancang sebagai repositori khusus untuk menyimpan seluruh \textit{audit record}. Selain berfungsi sebagai media penyimpanan, repositori ini juga mendukung mekanisme pengelolaan retensi sehingga \textit{audit record} dapat dipertahankan sebagai bukti digital selama periode retensi yang telah ditentukan.

Untuk mendukung proses audit maupun investigasi, auditor dapat mengakses \textit{Audit Trail Service} melalui mekanisme \textit{query}. Berdasarkan parameter yang diberikan, sistem akan mengambil \textit{audit record} yang sesuai dari repositori audit dan mengembalikannya kepada auditor. Penelitian ini hanya berfokus pada mekanisme penyediaan kembali (\textit{retrieval}) \textit{audit trail} sebagai bukti digital. Fungsi lanjutan seperti pemantauan (\textit{monitoring}), analisis aktivitas, maupun pemberian peringatan (\textit{alerting}) tidak termasuk ke dalam ruang lingkup sistem yang dikembangkan.

% \input{chapters/ta/chapter-3/rancangan-solusi/fungsionalitas.tex}

\section{Rancangan}

\subsection{Arsitektur Sistem Audit Trail}

Untuk mengatasi permasalahan yang telah diidentifikasi pada Analisis Masalah
--- yaitu tidak adanya sistem \textit{audit trail} yang terintegrasi, \textit{append-only},
dan mudah dianalisis --- penelitian ini mengusulkan penambahan satu komponen
baru pada arsitektur DecMed, yaitu \textbf{Audit Trail Service (ATS)}.
ATS dirancang sebagai layanan \textit{backend} yang independen terhadap
komponen DecMed yang telah ada (PRE Server, Gas Station Server, maupun
\textit{smart contract}), sehingga penambahan mekanisme audit trail tidak
mengubah alur bisnis maupun fungsionalitas DecMed yang sudah berjalan.
Arsitektur sistem audit trail dapat dilihat pada Gambar~\ref{fig:audit-arsitektur}.

ATS terdiri atas tiga modul utama:
\begin{itemize}
    \item \textbf{Audit Collector} -- menyediakan \textit{endpoint} penerimaan
    audit \textit{event} dari tiga sumber, yaitu Client Apps, PRE Server, dan
    Gas Station Server. Setiap sumber mengirimkan \textit{event} melalui
    permintaan terautentikasi ke Audit Collector.
    \item \textbf{Audit Processor} -- bertugas memvalidasi, menstandardisasi,
    dan melengkapi \textit{event} yang diterima menjadi \textit{audit record}
    yang siap disimpan, termasuk menghitung \textit{hash} dari isi \textit{record}
    sebagai bukti integritas awal.
    \item \textbf{Audit Storage Manager} -- bertugas mengunggah \textit{audit
    record} ke IPFS Cluster serta menginisialisasi transaksi bersponsor ke
    jaringan IOTA melalui Gas Station Server yang sudah ada, mengikuti pola
    penyimpanan yang sama dengan data RME.
    \item \textbf{Audit Query Index} -- basis data sekunder (non-otoritatif)
    yang menyimpan salinan metadata audit record untuk mendukung kebutuhan
    pencarian dan analisis secara efisien, tanpa mengubah sifat \textit{append-only}
    dari data yang tersimpan di IOTA dan IPFS.
\end{itemize}

ATS berkomunikasi dengan Gas Station Server yang telah ada pada arsitektur
DecMed untuk mensponsori transaksi audit trail ke jaringan IOTA, serta dengan
IPFS Cluster yang telah ada untuk menyimpan berkas \textit{audit record}.
Dengan demikian, ATS tidak menambah infrastruktur penyimpanan baru, melainkan
memanfaatkan kembali (\textit{reuse}) infrastruktur \textit{off-chain} dan
\textit{on-chain} yang sudah tersedia pada DecMed.

Selain itu, ditambahkan satu modul baru pada \textit{smart contract} DecMed,
yaitu \texttt{decmed::audit\_trail}, yang menyediakan \textit{method}
\texttt{create\_audit\_entry} untuk mencatat referensi \textit{audit record}
(CID beserta metadata terkait) pada penyimpanan \textit{on-chain}. Modul ini
hanya menyediakan operasi penciptaan entri baru dan secara sengaja tidak
menyediakan \textit{method} pengubahan maupun penghapusan, sehingga setiap
entri yang telah tercatat bersifat permanen (\textit{append-only}).

\subsection{Mekanisme Pembangkitan Audit Trail}

Audit \textit{event} dibangkitkan dari tiga sumber yang telah diidentifikasi
pada tahap analisis kebutuhan audit trail (III.1.3), yaitu:

\begin{enumerate}
    \item \textbf{Client Apps} -- membangkitkan \textit{event} atas aktivitas
    yang dilakukan langsung oleh pengguna pada perangkatnya, seperti
    autentikasi lokal (PIN/biometrik), registrasi akun, dan pemindaian kode QR
    delegasi akses.
    \item \textbf{PRE Server} -- membangkitkan \textit{event} atas setiap
    permintaan yang ditanganinya, seperti proses autentikasi (\textit{nonce},
    penerbitan token), pengambilan maupun perubahan data RME, dan re-enkripsi
    data.
    \item \textbf{Gas Station Server} -- membangkitkan \textit{event} atas
    setiap permintaan reservasi dan eksekusi transaksi bersponsor yang
    ditanganinya.
\end{enumerate}

Setiap sumber membangkitkan \textit{event} dalam format standar yang telah
ditentukan, mengacu pada kebutuhan isi \textit{audit record} sebagaimana
dianalisis pada III.1.3 dan selaras dengan kontrol AU-3 (\textit{Content of
Audit Records}) pada NIST SP 800-53 Rev. 5, yang mencakup atribut berikut:

\begin{xltabular}{\textwidth}{|>{\centering\arraybackslash}p{2.8cm}|X|}
    \hline
    \textbf{Atribut} & \textbf{Deskripsi} \\
    \hline
    \texttt{event\_type} & Jenis aktivitas yang dicatat (mis. \texttt{AUTH\_LOGIN},
    \texttt{ACCESS\_GRANT}, \texttt{ACCESS\_REVOKE}, \texttt{RECORD\_READ},
    \texttt{RECORD\_WRITE}, \texttt{KEY\_ISSUANCE}, \texttt{TX\_SPONSOR}, dsb.) \\
    \hline
    \texttt{source} & Sumber pembangkit \textit{event}, yaitu \texttt{CLIENT\_APP},
    \texttt{PRE\_SERVER}, atau \texttt{GAS\_STATION} \\
    \hline
    \texttt{actor\_ref} & Identitas atau referensi aktor yang terlibat (mis.
    \texttt{iota\_address} pasien/personel) \\
    \hline
    \texttt{timestamp} & Waktu terjadinya aktivitas \\
    \hline
    \texttt{outcome} & Status hasil aktivitas (\textit{success}/\textit{failure}) \\
    \hline
    \texttt{detail} & Data kontekstual tambahan spesifik \textit{event}
    (mis. \texttt{capability\_id}, \texttt{medical\_metadata\_index}) \\
    \hline
\end{xltabular}

Setelah \textit{event} dibangkitkan oleh sumber, \textit{event} tersebut
dikirimkan ke Audit Collector pada ATS melalui permintaan terautentikasi.
Audit Processor kemudian memvalidasi kelengkapan atribut, melengkapi
\texttt{audit\_id} unik, dan menghitung \texttt{record\_hash} dari
keseluruhan isi \textit{record} sebelum diteruskan ke tahap penyimpanan.

\subsection{Mekanisme Penyimpanan Audit Trail}

Mekanisme penyimpanan \textit{audit trail} mengikuti pola yang sama dengan
mekanisme penyimpanan data RME pada DecMed (III.2.2 pada dokumen rancangan
DecMed), yaitu memisahkan penyimpanan isi data (\textit{off-chain}) dari
penyimpanan referensi dan metadata (\textit{on-chain}).

Setiap \textit{audit record} yang telah distandardisasi oleh Audit Processor
diserialisasi menjadi satu berkas (satu \textit{event} menghasilkan satu
berkas) dan diunggah ke IPFS Cluster yang telah digunakan DecMed untuk
menyimpan data klinis RME. Proses unggah ini menghasilkan \textit{Content
Identifier} (CID) yang bersifat unik terhadap isi berkas tersebut; apabila
isi berkas berubah, CID yang dihasilkan juga akan berbeda, sehingga CID
berfungsi sebagai \textit{fingerprint} integritas \textit{audit record}.

Selanjutnya, Audit Storage Manager menginisialisasi transaksi bersponsor
melalui Gas Station Server untuk memanggil \texttt{create\_audit\_entry}
pada modul \texttt{decmed::audit\_trail}. Transaksi tersebut menyimpan CID
beserta metadata \textit{audit record} (\texttt{event\_type}, \texttt{source},
\texttt{actor\_ref}, \texttt{timestamp}, \texttt{record\_hash}) pada
penyimpanan \textit{on-chain} IOTA. Karena modul \texttt{decmed::audit\_trail}
hanya menyediakan operasi penciptaan entri baru, setiap aktivitas -- termasuk
pemberian dan pencabutan akses -- akan selalu tercatat sebagai entri baru
yang berdiri sendiri, alih-alih memperbarui entri yang sudah ada seperti pada
\textit{PatientAccessLog}. Pendekatan ini secara langsung menjawab kelemahan
\textit{immutability} yang diuraikan pada Analisis Masalah.

Dengan kombinasi CID (di IPFS) dan referensi \textit{on-chain} (di IOTA),
setiap \textit{audit record} memiliki dua lapis jaminan integritas: isi
\textit{record} tidak dapat diubah tanpa mengubah CID-nya, dan CID yang
telah dicatat pada IOTA tidak dapat diubah maupun dihapus setelah transaksi
final. Dokumen ini tidak membahas mekanisme \textit{batching} maupun
optimasi kinerja pengiriman, karena isu kinerja dan skalabilitas berada di
luar cakupan pembahasan penelitian ini.

\subsection{Mekanisme Pengelolaan Audit Trail}

Mekanisme pengelolaan \textit{audit trail} pada ATS mencakup tiga aspek,
yaitu pencarian/analisis, retensi, dan verifikasi integritas.

\subsubsection{Pencarian dan Analisis}
Data \textit{on-chain} dan \textit{off-chain} tetap menjadi sumber
kebenaran (\textit{source of truth}) yang bersifat \textit{immutable}, namun
tidak efisien untuk ditelusuri secara langsung. Oleh karena itu, setiap kali
sebuah \textit{audit record} berhasil disimpan, ATS turut menyimpan salinan
metadatanya (bukan isi lengkapnya) pada Audit Query Index, yaitu basis data
sekunder yang mendukung pencarian dan agregasi berdasarkan atribut seperti
rentang waktu, jenis \textit{event}, sumber, maupun aktor. Mekanisme ini
menjawab kelemahan \textit{transaction block} IOTA yang tidak memiliki
mekanisme \textit{query}, indeks, maupun agregasi log.

\subsubsection{Retensi}
Berkas \textit{audit record} pada IPFS Cluster di-\textit{pin} dengan
kebijakan yang sama dengan data klinis RME, sehingga tidak hilang akibat
mekanisme \textit{garbage collection} pada node IPFS yang tidak melakukan
\textit{pinning}. Referensi \textit{on-chain} pada IOTA tidak memiliki
mekanisme penghapusan pada level \textit{smart contract}, sehingga retensi
metadata audit trail bergantung pada kebijakan retensi ledger IOTA itu
sendiri di luar level aplikasi.

\subsubsection{Verifikasi Integritas}
Pihak yang berwenang (auditor/administrator) dapat memverifikasi keutuhan
sebuah \textit{audit record} secara independen dengan mengambil CID yang
tercatat pada IOTA, mengambil isi berkas dari IPFS menggunakan CID tersebut,
menghitung ulang \textit{hash} isi berkas, dan membandingkannya dengan
\texttt{record\_hash} yang tercatat \textit{on-chain}. Kecocokan hasil
verifikasi menandakan bahwa \textit{audit record} belum mengalami perubahan
sejak pertama kali dicatat.

\subsection{Fungsionalitas Sistem}

Subbab ini menjelaskan fungsionalitas sistem audit trail secara rinci,
mencakup deskripsi alur bagi setiap kebutuhan fungsional yang telah
diidentifikasi pada III.1.3, kecuali kebutuhan fungsional yang bersifat
sederhana (mis. validasi format masukan) yang tidak memerlukan penjelasan
alur tersendiri. Setiap mekanisme direferensikan ke diagram sekuens terkait,
yang akan dilengkapi pada bagian berikutnya.

\subsubsection{Mekanisme Pembangkitan dan Pengiriman Audit Record dari Client Apps}
Mekanisme ini dapat dilihat pada Gambar~\ref{fig:audit-generate-client}.
Ketika pengguna melakukan aktivitas yang tercakup dalam cakupan audit pada
Client Apps (mis. autentikasi lokal PIN/biometrik, registrasi akun, atau
pemindaian kode QR delegasi akses), Client Apps menyusun \textit{event}
dengan atribut standar (\texttt{event\_type}, \texttt{source}=\texttt{CLIENT\_APP},
\texttt{actor\_ref}, \texttt{timestamp}, \texttt{outcome}, \texttt{detail}) dan
mengirimkannya ke \textit{endpoint} Audit Collector pada ATS melalui
permintaan terautentikasi. Audit Collector menerima \textit{event} tersebut
dan meneruskannya ke Audit Processor untuk tahap pemrosesan lebih lanjut,
sebagaimana dijelaskan pada mekanisme pemrosesan dan penyimpanan.

\subsubsection{Mekanisme Pembangkitan dan Pengiriman Audit Record dari PRE Server}
Mekanisme ini dapat dilihat pada Gambar~\ref{fig:audit-generate-pre}.
Setiap kali PRE Server menangani permintaan pada salah satu \textit{endpoint}-nya
(\texttt{/nonce}, \texttt{/keys}, \texttt{/medical-record}, \texttt{/administrative},
dan sejenisnya), PRE Server menyusun \textit{event} dengan
\texttt{source}=\texttt{PRE\_SERVER} yang merepresentasikan aktivitas
autentikasi, penerbitan token, maupun akses/perubahan data RME yang telah
ditangani, kemudian mengirimkan \textit{event} tersebut ke Audit Collector
pada ATS. Pengiriman \textit{event} dilakukan secara asinkron terhadap alur
utama permintaan pengguna, sehingga tidak menghambat waktu respons PRE
Server terhadap Client Apps.

\subsubsection{Mekanisme Pembangkitan dan Pengiriman Audit Record dari Gas Station Server}
Mekanisme ini dapat dilihat pada Gambar~\ref{fig:audit-generate-gas}.
Setiap kali Gas Station Server menerima permintaan reservasi maupun eksekusi
\textit{gas object} untuk mensponsori suatu transaksi, Gas Station Server
menyusun \textit{event} dengan \texttt{source}=\texttt{GAS\_STATION} yang
mencatat transaksi yang disponsorinya (termasuk transaksi
\texttt{create\_audit\_entry} milik ATS sendiri), kemudian mengirimkan
\textit{event} tersebut ke Audit Collector pada ATS.

\subsubsection{Mekanisme Pemrosesan dan Penyimpanan Audit Record}
Mekanisme ini dapat dilihat pada Gambar~\ref{fig:audit-process-store}.
Setelah \textit{event} diterima oleh Audit Collector dari salah satu dari
ketiga sumber di atas, Audit Processor memvalidasi kelengkapan atribut
\textit{event}, melengkapi \texttt{audit\_id} unik apabila belum tersedia,
dan menghitung \texttt{record\_hash} dari keseluruhan isi \textit{record}.
\textit{Record} yang telah diproses diserahkan kepada Audit Storage Manager,
yang menyerialisasikannya menjadi satu berkas (satu \textit{event}
menghasilkan satu berkas) dan mengunggah berkas tersebut ke IPFS Cluster,
menghasilkan CID. Audit Storage Manager selanjutnya menginisialisasi
transaksi bersponsor yang memanggil \texttt{create\_audit\_entry} pada modul
\texttt{decmed::audit\_trail} melalui Gas Station Server. Setelah Gas
Station Server mengeksekusi transaksi tersebut dan dikonfirmasi oleh
jaringan IOTA, CID beserta metadata \textit{audit record} tercatat secara
permanen pada penyimpanan \textit{on-chain}. ATS kemudian memperbarui Audit
Query Index dengan metadata yang sama untuk mendukung kebutuhan pencarian.

\subsubsection{Mekanisme Pencarian Audit Trail}
Mekanisme ini dapat dilihat pada Gambar~\ref{fig:audit-query}.
Pihak berwenang (Admin Fasyankes/Kementerian Kesehatan) memasukkan kriteria
pencarian, seperti rentang waktu, jenis \textit{event}, sumber, maupun
aktor, melalui antarmuka yang tersedia pada Client Fasyankes/Kementerian.
Permintaan pencarian diteruskan ke ATS, yang menelusuri Audit Query Index
untuk memperoleh daftar \textit{audit record} yang sesuai kriteria. ATS
mengembalikan metadata beserta referensi CID dari setiap \textit{record}
yang ditemukan untuk ditampilkan kepada pengguna.

\subsubsection{Mekanisme Verifikasi Integritas Audit Record}
Mekanisme ini dapat dilihat pada Gambar~\ref{fig:audit-verify}.
Pengguna memilih satu \textit{audit record} hasil pencarian untuk
diverifikasi keutuhannya. ATS mengambil CID \textit{record} tersebut dari
penyimpanan \textit{on-chain} IOTA, mengambil isi berkas dari IPFS
menggunakan CID tersebut, menghitung ulang \textit{hash} isi berkas, dan
membandingkannya dengan \texttt{record\_hash} yang tercatat \textit{on-chain}.
Hasil perbandingan (\textit{match}/\textit{mismatch}) dikembalikan dan
ditampilkan kepada pengguna sebagai bukti keutuhan \textit{audit record}
sejak pertama kali dicatat.